% EPL master thesis covers template
\documentclass{template/EPL-master-thesis-covers-EN}
\usepackage{fancyhdr}
\usepackage{tabularx}
\usepackage{geometry}
\usepackage{rotating}
\usepackage{minitoc}
\usepackage{verbatim}
\usepackage{pdfpages}
\usepackage{float}
\usepackage{listings}
\usepackage{graphics}
\usepackage{listings}
\usepackage{geometry}
\usepackage{mathtools}
\usepackage{amsmath}
\usepackage{amssymb}
\usepackage{newtxtext,newtxmath}
\usepackage{siunitx}
\usepackage{float}
\usepackage[inkscapeformat=png]{svg}
\usepackage[backend=bibtex, style=ieee]{biblatex}
\usepackage{csquotes}
\usepackage{hyperref}


% Please fill in the following boxes
% Title of the thesis
\title{Direct device access from the smartNIC towards datacenters disaggregation}

% Subtitle - remove this line if not applicable
\subtitle{Optional subtitle}

% Name of the student author(s)
\author{Nicolas \textsc{Jeanmenne}}
% \secondauthor{Firstname \textsc{Lastname}}		% remove if not applicable
%\thirdauthor{Firstname \textsc{Lastname}}			% remove if not applicable

% Official title of the master degree (copy/paste from list below)
% Master [120] in Biomedical Engineering
% Master [120] in Chemical and Materials Engineering
% Master [120] in Civil Engineering
% Master [120] in Computer Science
% Master [120] in Computer Science and Engineering
% Master [120] in Cybersecurity
% Master [120] in Data Sciences Engineering
% Master [120] in Data Science: Information technology
% Master [120] in Electrical Engineering
% Master [120] in Electro-mechanical Engineering
% Master [120] in Mathematical Engineering
% Master [120] in Mechanical Engineering
% Master [120] in Physical Engineering
% Master [60] in Computer Science
% Specialised master in nanotechnologies
% Specialised master in nuclear engineering
\degreetitle{Master [120] in Computer Science}

% Name of the supervisor(s)
\supervisor{Tom \textsc{Barbette}}
% \secondsupervisor{Firstname \textsc{Lastname}}		% remove if not applicable
%\thirdsupervisor{Firstname \textsc{Lastname}}		% remove if not applicable

% Name of the reader(s)
\readerone{Firstname \textsc{Lastname}}
\readertwo{Firstname \textsc{Lastname}}			% remove if not applicable
\readerthree{Firstname \textsc{Lastname}}			% remove if not applicable
%\readerfour{Firstname \textsc{Lastname}}			% remove if not applicable
%\readerfive{Firstname \textsc{Lastname}}			% remove if not applicable

% Academic year (update if necessary)
\years{2025--2026}

\newcommand{\authcite}[1]{\citeauthor{#1} \cite{#1}}

\bibliography{biblio.bib}

% Document
\begin{document}
  % Front cover page
  \maketitle


\chapter*{Abstract}

  TO-DO: complete abstract

  \textit{This paper is draft aiming at analyzing state-of-the-art solution for DDA from the smartNIC towards datacenter disaggregation
  and is not made to be submitted anywhere}

\chapter{State-of-the-art}

\section{Introduction}

For decades, servers in datacenters were designed under a monolithic model, meaning that all hardware resources are physically
attached to the same frame. This model is now a bottleneck for performance, network throughput and to efficiently consume resources.
In order to solve the monolithic challenge, a new architectural design has been proposed : disaggregation.

\section{Definitions}

TO-DO : rewrite without subsections

\subsection{Disaggregation}

Disaggregation is a new model where hardware resources such as computing power, memory, accelerators, storage, ... are divided
into pools (or nodes) interconnected into each other via a high-speed network. This allows resources to be allocated on demand,
with better utilization and availability. There is mainly two levels of disaggregation.

Partial disaggregation separates storage components from CPU and memory which remain integrated together. This level of
disaggregation has already been widely implemented since early 2020s.

Fully disaggregation regroups same type of resource into clusters which are interconnected over a network to allow
communication and data exchange between them. \cite{lin_disaggregated_2020}

\subsection{RDMA}

Remote Direct Memory Access (RDMA) is a protocol that allows accessing memory from one computer into another computer without
involving operating system \cite{recio_remote_2007}. This protocol offers high throughput and low latency, and will be referred from many articles
covering this topic.

\subsection{Zero-copy}

Zero-copy enable the possibility to transfer data from one device to another without requiring CPU to copy the data. This
avoids redundant copies, reducing CPU usage.

\subsection{SPDK}

TO-DO : complete

\subsection{NVMe-oF}

NVMe over-Fabric (NVMe-oF) is an extension of NVMe command set. It allows sending these commands over networked fabrics
\cite{noauthor_nvme_2021, noauthor_what_nodate, guz_nvme-over-fabrics_2017}

\section{Architectural challenges}

Challenges whose origin is architectural design can be divided into two classes :


\begin{itemize}
  \item Problems due to traditional monolithic architectural server. This type of problem is the core reason of why datacenters
    are moving from monolithic architecture from disaggregated one.
  \item New challenges brought by disaggregation. Disaggregated computing is a relatively new design in system and datacenter
    who breaks many of assumption previously made in the research. Researchers and datacenter operators need to find
    creative solutions to outcome these challenges and enjoy the benefits of disaggregation.
\end{itemize}

\subsection{Monolithic challenges}

The main drawback with monolithic architecture (MA) is \textbf{resource stranding}. When a server exhausts one of its resource,
for example memory, all other resources cannot be utilized and result in a waste of resources.

Another non-negligible issue concerning MA is failure tolerance and availability. When a piece of hardware fails this may cause
the entire server to fails and become unavailable.

Finally, replacing or upgrading hardware in MA occurs maintenance that can be costly.

\subsection{Disaggregation challenges}

One of the main challenge in disaggregated architecture (DA) is that network bandwidth becomes a central point of communication between
clusters; CPU-storage communication needs to be done in a few milliseconds while CPU-memory needs to be done in a few nanoseconds, usually
100ns \cite{lin_disaggregated_2020}.

Existing network stacks face significant trade-offs in terms of flexibility and performance. For example, RDMA and RDMA
over Convergent Ethernet (RoCE) are high-performance protocols, but they lack flexibility because they are tightly coupled
to specific hardware, making them difficult to adapt. In addition to limited flexibility, RDMA suffers from issues like
head-of-line blocking, congestion, and vendor lock-in, due to hardware requirements that aren't met by commodity equipment.
On the other hand, network stacks like Linux TCP offer greater flexibility but at the cost of poor performance and
inefficient CPU utilization \cite{skiadopoulos_high-throughput_2024}.

\section{State-of-the-art designs and solutions} \label{sec:sota}

Some hardware pieces are evolving, like storage moving from hard disk drive to NVMe SSD and RAM becoming less expensive.
This allows researchers and enterprises to discover several new ways to achieve high performance computing. In this section
we'll review main contributions from the literature.

\subsection{Zero-copy data path}

\authcite{skiadopoulos_high-throughput_2024} designed a new prototype called \textit{ZeroNIC} which can zero-copy data between
NIC and accelerators. Its main advantage is that is it agnostic of the accelerator type and can run on CPU, GPU, FPGA, ...
NIC hardware split header and payload, transferring the latter directly to the receiver applications buffers without intermediates
copies.

Concerning NVMe storage, \authcite{sun_scalio_2025} focused on NVMe-oF target offloading which extends zero-copy mechanism,
transferring data to the DPU's host channel adapter (HCA) via peer-to-peer PCI communication (P2PDMA) and thus creating a new
kind of zero-copy data path.

\subsection{Batched write}

Batched write mechanism retains write operations until an arbitrary threshold is achieved, then all these operations are run
consecutively while a group commit ensures atomicity of operations. The goal is to improve throughput and reduce the accelerator
CPU load. While it is not technically a contribution brought by disaggregation, its evolution into such systems is still relevant
to be mentioned.

\subsection{Split stacks}

Split stacks are networking stacks divided in several planes, each one optimized for a specific task. IO-TCP from
\authcite{kim_rearchitecting_2023} separate the traditional Linux TCP stack into a stack running on the host CPU, performing
control operations, and thus called the control plane. Such operations are complex and benefits from features of the CISC CPU
(for example branch prediction, vectorized instructions, ...). The host stack extends the mTCP \cite{jeong_mtcp_2014} API with methods to offload
file opening, writing, and closing, such that applications have the most minimal changes to do in order to run itself on IO-TCP.

The data plane handles all operations related to data packet preparation and transfer. These operations are simpler in comparison
of the control plane and then fit better the ARM cores of the smartNIC. They also benefit from being stateless which suits the
usage of the smartNIC. File and disk I/O are offloaded to otherwise it would interfere with the control plane on host CPU.
Data plane operations are virtually performed and enable a mechanism of SEND / ECHO which translate in multiple MTU packets
transfer.

Many difficulties due to the split introduce themselves, like RTT measurement, retransmission, and error handling, but are
essentially resolved with the set of commands provided by the implementation.

\section{Unresolved challenges}

\paragraph{}
There are some challenges still left open upon reading papers in \textbf{section \ref{sec:sota}}. This section aims to
list them, with the goal to define a clear approach for working on the master's thesis as well as a precise problem to solve.

IO-TCP \cite{kim_rearchitecting_2023} lacks of dynamic synchronization. This because authors assumes that files don't change
during the content delivery service. In their solution, both host and smartNIC share the same filesystem, but any change
on the host-side FS isn't propagated to the smartNIC stack.
% -> Distributed FS ?

Another problem is that file reading on NVMe disks is done with direct I/O on a regular ext4 file system. SPDK reaches the best
performance, but support for file system isn't very developed for now.

Next, with a smaller amount of connection (typically < 4), Linux TCP stack outperforms IO-TCP by over 1.5 times.

\paragraph{}
Batched writes used in Scalio \cite{sun_scalio_2025} to improve throughput of write operations come at the cost of latency
increase. This can lead to non-negligible issues in environment where low latency is critical.

\section{Objectives}

This paper will focus on fully disaggregated datacenter, especially on interaction between NVMe storage pools and smartNIC.
One of our objectives is to allow SmartNICs to communicate directly with the storage without involving any compute pool.
% TO-DO:  give a more precise definition

\printbibliography
 
  % Back cover page
  \backcoverpage

\end{document}
